\chapter{Alat: računske radnje koje ćeš stalno ponavljati}
\label{ch:alat}

Kroz sva četiri zadatka ponavlja se svega desetak računskih radnji. Ovdje su
sve izvedene do kraja, na brojevima iz seminara, tako da ih poslije ne moraš
iznova izvoditi --- samo se pozoveš na ovo poglavlje.

\section{Množenje matrice i vektora}
\label{sec:mnozenje}

\begin{formula}{Matrica puta vektor}
Za $M \in \R^{m \times n}$ i $x \in \R^{n}$ rezultat $Mx$ je vektor iz $\R^{m}$
kojemu je $i$-ta komponenta skalarni produkt $i$-tog \emph{retka} matrice s
vektorom $x$:
\begin{equation}
(Mx)_i \;=\; \sum_{k=1}^{n} M_{ik}\, x_k .
\end{equation}
\end{formula}

Pravilo za pamćenje: \emph{redak puta stupac}. Broj stupaca lijeve matrice mora
biti jednak broju redaka desne, inače množenje nije definirano.

Na primjeru iz Zadatka 1, s $G = \begin{bsmallmatrix} -1 & -1 \\ 1 & -1 \end{bsmallmatrix}$
i $x = \begin{bsmallmatrix} x_1 \\ x_2 \end{bsmallmatrix}$:
\begin{equation}
Gx = \begin{bmatrix} -1 & -1 \\ 1 & -1 \end{bmatrix}
     \begin{bmatrix} x_1 \\ x_2 \end{bmatrix}
   = \begin{bmatrix} -x_1 - x_2 \\ x_1 - x_2 \end{bmatrix}.
\end{equation}
Prvi redak $\begin{bsmallmatrix} -1 & -1\end{bsmallmatrix}$ pomnožen s $x$ daje
$-x_1 - x_2$; drugi redak daje $x_1 - x_2$.

\section{Kvadratna forma i zapis funkcije cilja u matričnom obliku}
\label{sec:kvadratna-forma}

Ovo je prvi korak svakog QP zadatka i mjesto gdje se najčešće promaši faktor.

\begin{pojam}{Kvadratna forma}
Za simetričnu matricu $Q \in \Sym{n}$ izraz
\begin{equation}
x\trans Q x \;=\; \sum_{i=1}^{n}\sum_{j=1}^{n} Q_{ij}\, x_i x_j
\end{equation}
zove se \textbf{kvadratna forma}. Rezultat je jedan broj. Za $n = 2$ raspisano:
\begin{equation}
x\trans Q x = Q_{11}x_1^2 + 2Q_{12}x_1x_2 + Q_{22}x_2^2 ,
\label{eq:kv-forma-2d}
\end{equation}
gdje se mješoviti član pojavljuje \emph{dvaput} jer je $Q_{12} = Q_{21}$.
\end{pojam}

\begin{korak}[Kako iz zadane funkcije pročitati $Q$ i $c$]
Tražimo zapis $f(x) = \tfrac{1}{2} x\trans Q x + c\trans x$. Postupak:
\begin{enumerate}[leftmargin=*]
  \item \textbf{Linearni članovi} idu izravno u $c$: koeficijent uz $x_i$ je $c_i$.
  \item \textbf{Čisto kvadratni članovi}: koeficijent uz $x_i^2$ jednak je
        $\tfrac{1}{2}Q_{ii}$, dakle $Q_{ii}$ je \emph{dvostruki} koeficijent uz $x_i^2$.
  \item \textbf{Mješoviti član}: koeficijent uz $x_i x_j$ jednak je
        $\tfrac{1}{2}\cdot 2 Q_{ij} = Q_{ij}$, dakle $Q_{ij}$ je \emph{isti}
        koeficijent uz $x_ix_j$, podijeljen između $Q_{ij}$ i $Q_{ji}$ tako da
        matrica ostane simetrična.
\end{enumerate}
\end{korak}

Primijenjeno na funkciju cilja iz Zadatka 1,
$f(x) = x_1^2 + 2x_2^2 - 0{,}3\,x_1x_2 + 2x_1 - 3x_2$:
\begin{align}
Q_{11} &= 2 \cdot 1 = 2, &
Q_{22} &= 2 \cdot 2 = 4, &
Q_{12} = Q_{21} &= -0{,}3, \\
c_1 &= 2, &
c_2 &= -3. & &
\end{align}
Dakle
\begin{equation}
Q = \begin{bmatrix} 2 & -0{,}3 \\ -0{,}3 & 4 \end{bmatrix}, \qquad
c = \begin{bmatrix} 2 \\ -3 \end{bmatrix}.
\label{eq:alat-Q-c}
\end{equation}

\textbf{Provjera unatrag} --- uvijek je napravi, traje deset sekundi:
\begin{equation}
\tfrac{1}{2}x\trans Q x = \tfrac{1}{2}\bigl(2x_1^2 + 2(-0{,}3)x_1x_2 + 4x_2^2\bigr)
= x_1^2 - 0{,}3\,x_1x_2 + 2x_2^2 . \quad\checkmark
\end{equation}

\begin{zamka}[Zamka: faktor $\tfrac{1}{2}$ i dijeljenje mješovitog člana]
Dvije greške koje se rade zajedno i međusobno se ,,ponište`` pa ostanu
neprimijećene do trenutka kad rezultat ispadne kriv:
\begin{itemize}[nosep,leftmargin=*]
  \item Zaboravi se da standardni oblik ima $\tfrac{1}{2}$, pa se napiše
        $Q_{11} = 1$ umjesto $2$. Tada Hesseova matrica ispadne
        $\begin{bsmallmatrix} 1 & \cdot \\ \cdot & \cdot \end{bsmallmatrix}$,
        a Hesseova matrica \emph{jest} $Q$ --- ne polovina od $Q$.
  \item Mješoviti član $-0{,}3\,x_1x_2$ se podijeli na pola, pa se napiše
        $Q_{12} = -0{,}15$. Pogrešno: pola je već uračunato u faktor
        $\tfrac{1}{2}$ ispred cijele forme, a $Q_{12}$ i $Q_{21}$ zajedno daju
        $2Q_{12}$ prema \eqref{eq:kv-forma-2d}.
\end{itemize}
Lijek je uvijek isti: raspiši $\tfrac{1}{2}x\trans Q x$ natrag i usporedi sa
zadanom funkcijom.
\end{zamka}

\section{Gradijent i Hesseova matrica}
\label{sec:gradijent-hessian}

\begin{izvor}{Predavanje 2 --- Definicije, klasifikacija, konveksnost, slajd 33/71 (PDF str.~34)}
\slajd[0.86]{p02-gradijent-def}

\textbf{Transkripcija.} Neka je $f: \R^n \to \R$ diferencijabilna funkcija na
svojoj domeni. Gradijent je definiran kao vektor
\[
\grad f(x) \coloneqq \bigl[Df(x)\bigr]\trans =
\begin{bmatrix} \dfrac{\partial f(x)}{\partial x_1} \\[6pt] \vdots \\[2pt]
                \dfrac{\partial f(x)}{\partial x_n} \end{bmatrix}.
\]
Afina funkcija $g(x) = f(x_0) + \grad f(x_0)\trans (x - x_0)$ je aproksimacija
prvog reda funkcije $f$ u točki $x_0$.

\textbf{Što to znači:} $Df(x)$ je Jacobijeva matrica, a za skalarnu funkciju ona
je \emph{redak}. Gradijent je njezin transponat, dakle \emph{stupac}. Zato u
svim formulama gradijent stoji kao stupčani vektor, a njegov skalarni produkt s
nekim smjerom $d$ piše se $\grad f(x)\trans d$.
\end{izvor}

\begin{formula}{Gradijent i Hessian kvadratne funkcije}
Za $f(x) = \tfrac{1}{2}x\trans Q x + c\trans x$ sa simetričnim $Q$ vrijedi
\begin{equation}
\grad f(x) = Qx + c, \qquad \grad^2 f(x) = Q .
\label{eq:grad-hess-kvadratne}
\end{equation}
Hesseova matrica je \emph{konstantna} --- ne ovisi o $x$.
\end{formula}

To je razlog zašto se kod QP-a konveksnost provjerava jednom zauvijek, a ne u
svakoj točki posebno.

\begin{korak}[Provjera formule \eqref{eq:grad-hess-kvadratne} deriviranjem ,,rukom``]
Za funkciju iz Zadatka 1:
\begin{enumerate}[leftmargin=*]
  \item $\dfrac{\partial f}{\partial x_1}
        = \dfrac{\partial}{\partial x_1}\bigl(x_1^2 + 2x_2^2 - 0{,}3x_1x_2 + 2x_1 - 3x_2\bigr)
        = 2x_1 - 0{,}3x_2 + 2.$
  \item $\dfrac{\partial f}{\partial x_2} = 4x_2 - 0{,}3x_1 - 3.$
  \item Složeno u vektor:
        $\grad f(x) = \begin{bsmallmatrix} 2x_1 - 0{,}3x_2 + 2 \\ -0{,}3x_1 + 4x_2 - 3 \end{bsmallmatrix}$.
  \item Isto to preko \eqref{eq:grad-hess-kvadratne}:
        $Qx + c = \begin{bsmallmatrix} 2 & -0{,}3 \\ -0{,}3 & 4\end{bsmallmatrix}
                  \begin{bsmallmatrix} x_1 \\ x_2 \end{bsmallmatrix}
                + \begin{bsmallmatrix} 2 \\ -3 \end{bsmallmatrix}
                = \begin{bsmallmatrix} 2x_1 - 0{,}3x_2 + 2 \\ -0{,}3x_1 + 4x_2 - 3\end{bsmallmatrix}$. \checkmark
  \item Deriviraj još jednom: $\dfrac{\partial^2 f}{\partial x_1^2} = 2$,
        $\dfrac{\partial^2 f}{\partial x_1 \partial x_2} = -0{,}3$,
        $\dfrac{\partial^2 f}{\partial x_2^2} = 4$, pa je
        $\grad^2 f = \begin{bsmallmatrix} 2 & -0{,}3 \\ -0{,}3 & 4\end{bsmallmatrix} = Q$. \checkmark
\end{enumerate}
\end{korak}

Za afinu funkciju $g(x) = a\trans x + \beta$ gradijent je konstantan,
$\grad g(x) = a$, a Hesseova matrica je nul-matrica. Zato su gradijenti
ograničenja u Zadatku 1 fiksne strelice --- ne mijenjaju se od točke do točke.

\section{Svojstvene vrijednosti matrice $2\times 2$}
\label{sec:svojstvene}

\begin{formula}{Karakteristična jednadžba}
Svojstvene vrijednosti matrice $M$ su rješenja jednadžbe
\begin{equation}
\det(\lambda I - M) = 0 .
\end{equation}
Za $M = \begin{bsmallmatrix} a & b \\ b & d \end{bsmallmatrix}$ (simetričnu) to je
kvadratna jednadžba
\begin{equation}
\lambda^2 - (a+d)\,\lambda + (ad - b^2) = 0,
\label{eq:karakteristicna-2x2}
\end{equation}
gdje je $a + d = \tr(M)$ i $ad - b^2 = \det(M)$.
\end{formula}

Za $Q$ iz \eqref{eq:alat-Q-c}: $\tr(Q) = 2 + 4 = 6$, $\det(Q) = 2\cdot 4 - (-0{,}3)^2 = 7{,}91$, pa je
\begin{equation}
\lambda^2 - 6\lambda + 7{,}91 = 0
\;\Longrightarrow\;
\lambda_{1,2} = \frac{6 \pm \sqrt{36 - 31{,}64}}{2} = \frac{6 \pm \sqrt{4{,}36}}{2},
\end{equation}
\begin{equation}
\lambda_1 = 1{,}955969, \qquad \lambda_2 = 4{,}044031 .
\label{eq:alat-svojstvene-Q}
\end{equation}
Obje su strogo pozitivne.

\begin{zamka}[Zamka: kad su svojstvene vrijednosti stvarno realne]
Formula \eqref{eq:karakteristicna-2x2} može dati kompleksna rješenja ako je
diskriminanta negativna. Za \emph{simetrične} matrice to se ne može dogoditi ---
diskriminanta je $(a-d)^2 + 4b^2 \geq 0$. Zato kod provjere konveksnosti nikad
ne završiš s kompleksnim brojevima. Kod \emph{nesimetričnih} matrica, kao što su
$A_{\text{parno}}$ i $A_{\text{neparno}}$ u Zadatku 4, kompleksne svojstvene
vrijednosti su posve normalne i baš one daju oscilatorno ponašanje.
\end{zamka}

\section{Provjera definitnosti: dva puta do istog odgovora}
\label{sec:definitnost}

\begin{izvor}{Predavanje 2 --- Definicije, klasifikacija, konveksnost, slajd 49/71 (PDF str.~50)}
\slajd[0.86]{p02-psd-determinante}

\textbf{Transkripcija.} Matrica $M$ je \textbf{pozitivno definitna} ako i samo ako
su sve determinante u nizu vodećih glavnih minora pozitivne:
\[
D_1 = m_{11}, \quad
D_2 = \begin{vmatrix} m_{11} & m_{12} \\ m_{21} & m_{22} \end{vmatrix}, \quad
\dots, \quad
D_n = \det(M).
\]
$M$ je \textbf{negativno definitna} ako i samo ako predznaci alterniraju počevši
od negativnog: $D_1 < 0$, $D_2 > 0$, $D_3 < 0$, $D_4 > 0, \dots$

\textbf{Kako se to koristi:} ovaj kriterij ne traži rješavanje karakteristične
jednadžbe --- samo determinante. Za matrice $2\times2$ to su dva broja i gotov si
u dva reda.
\end{izvor}

Za $Q$ iz \eqref{eq:alat-Q-c}: $D_1 = 2 > 0$ i $D_2 = \det Q = 7{,}91 > 0$, pa je
$Q \pd 0$ --- isti zaključak kao iz \eqref{eq:alat-svojstvene-Q}, uz upola manje
računa.

\begin{izgovor}
,,Pozitivnu definitnost provjerio sam na dva neovisna načina: preko svojstvenih
vrijednosti, koje su $1{,}956$ i $4{,}044$, dakle obje strogo pozitivne, i preko
vodećih glavnih minora, $D_1 = 2$ i $D_2 = 7{,}91$, oba pozitivna. Oba kriterija
daju isti zaključak, matrica je pozitivno definitna.``
\end{izgovor}

\section{Inverz matrice $2\times2$ i rješavanje malog sustava}
\label{sec:inverz}

\begin{formula}{Inverz matrice $2\times2$}
\begin{equation}
M = \begin{bmatrix} a & b \\ c & d \end{bmatrix}
\;\Longrightarrow\;
M^{-1} = \frac{1}{ad - bc} \begin{bmatrix} d & -b \\ -c & a \end{bmatrix},
\qquad ad - bc \neq 0 .
\end{equation}
Pravilo: zamijeni dijagonalu, promijeni predznak vandijagonalnima, podijeli s
determinantom.
\end{formula}

Primijenjeno na $Q$: $\det Q = 7{,}91 = \tfrac{791}{100}$, pa je
\begin{equation}
Q^{-1} = \frac{100}{791} \begin{bmatrix} 4 & 0{,}3 \\ 0{,}3 & 2 \end{bmatrix}.
\label{eq:alat-Q-inv}
\end{equation}

Ovaj se inverz koristi na dva mjesta u Zadatku 1: kod nalaženja neograničenog
stacionarnog rješenja $x^{0} = -Q^{-1}c$ i kod izvoda dualne funkcije.

\begin{zamka}[Zamka: računanje u decimalama umjesto u razlomcima]
$\det Q = 7{,}91$ je egzaktno $791/100$. Ako od početka radiš u razlomcima,
rješenje Zadatka 1 ispada kao egzaktan razlomak $\xopt = (190/33,\; 140/33)$.
Ako radiš u decimalama i zaokružuješ, dobiješ $5{,}7576$ i $4{,}2424$ pa na
obrani ne možeš pokazati da je rezultat točan, nego samo da je približan. Uvijek
vodi egzaktan trag, a decimale piši uz njega u zagradi.
\end{zamka}

\section{Pravac, poluprostor i kako ograničenje pretvoriti u crtež}
\label{sec:poluprostor}

\begin{izvor}{Predavanje 2 --- Definicije, klasifikacija, konveksnost, slajd 65/71 (PDF str.~66)}
\slajd[0.86]{p02-poluprostor}

\textbf{Transkripcija.} Poluprostor je skup oblika
\[
H \coloneqq \{\, x \in \R^n \mid a\trans x \leq b \,\}
\]
za zadani $a \in \R^n$, $a \neq 0$, i $b \in \R$.

\textbf{Što to znači u ravnini:} jednadžba $a\trans x = b$ opisuje pravac. Taj
pravac dijeli ravninu na dvije polovice; nejednakost $a\trans x \leq b$ bira
jednu od njih, zajedno sa samim pravcem. Vektor $a$ je okomit na pravac i
pokazuje u polovicu koja je \emph{isključena} (u smjeru rasta $a\trans x$).
\end{izvor}

\begin{korak}[Od nejednakosti do crteža, u četiri poteza]
Za ograničenje $g_1(x) = -x_1 - x_2 + 10 \leq 0$ iz Zadatka 1:
\begin{enumerate}[leftmargin=*]
  \item \textbf{Rub}: stavi znak jednakosti, $-x_1 - x_2 + 10 = 0$, tj.\
        $x_1 + x_2 = 10$.
  \item \textbf{Dvije točke}: uvrsti $x_1 = 0 \Rightarrow x_2 = 10$, pa
        $x_2 = 0 \Rightarrow x_1 = 10$. Kroz $(0,10)$ i $(10,0)$ povuci pravac.
  \item \textbf{Koja strana}: uvrsti probnu točku, najlakše ishodište
        $x = (0,0)$: $g_1(0,0) = 10 > 0$. Ishodište \emph{ne} zadovoljava
        ograničenje, dakle dopustiva je suprotna strana --- ona iznad pravca.
  \item \textbf{Strelica gradijenta}: $\grad g_1 = \begin{bsmallmatrix} -1 \\ -1\end{bsmallmatrix}$;
        nacrtaj je iz točke na pravcu. Pokazuje okomito na pravac, prema
        \emph{nedopustivoj} strani.
\end{enumerate}
\end{korak}

Isti postupak za $g_2(x) = x_1 - x_2 - 3 \leq 0$: rub je $x_1 - x_2 = 3$, prolazi
kroz $(3,0)$ i $(0,-3)$; ishodište daje $g_2(0,0) = -3 < 0$, dakle ishodište
\emph{jest} dopustivo, pa je dopustiva strana ona koja sadrži ishodište;
gradijent je $\begin{bsmallmatrix} 1 \\ -1\end{bsmallmatrix}$.

\begin{izvor}{Predavanje 2 --- Definicije, klasifikacija, konveksnost, slajd 58/71 (PDF str.~59)}
\slajd[0.86]{p02-presjek-konveksnih}

\textbf{Transkripcija.} Presjek konveksnih skupova je konveksan skup. Posljedica
za karakterizaciju dopustivog prostora $\feas$: ako je
$\feas_i = \{x \mid g_i(x) \leq 0\}$ ili $\feas_i = \{x \mid h_i(x) = 0\}$, i ako
su svi $\feas_i$, $i = 1, \dots, m$, konveksni, tada će i dopustivi skup
$\feas = \feas_1 \cap \feas_2 \cap \dots \cap \feas_m$ biti konveksan.

\textbf{Zašto to trebamo:} ovo je jedini korak koji nedostaje da bi se iz ,,svako
ograničenje je afino`` zaključilo ,,dopustivi skup je konveksan``. Bez njega
tvrdnja o konveksnosti problema nije dokazana nego samo pretpostavljena.
\end{izvor}

\section{Popis: što provjeriti prije nego kreneš računati}
\label{sec:kontrolna-lista}

\begin{predznanje}
Prije nego u bilo kojem zadatku napišeš prvi red računa:
\begin{enumerate}[nosep,leftmargin=*]
  \item Jesu li sva ograničenja svedena na oblik ,,nešto $\leq 0$``? Ako je
        zadano $a\trans x \leq b$, prebaci: $a\trans x - b \leq 0$.
  \item Je li $Q$ simetrična? Ako nije, zamijeni je sa $\tfrac{1}{2}(Q + Q\trans)$ ---
        kvadratna forma se time ne mijenja.
  \item Je li faktor $\tfrac{1}{2}$ uračunat u $Q$?
  \item Znaš li dimenzije svakog simbola koji si napisao? ($x \in \R^2$,
        $\mu \in \R^2$, $G \in \R^{2\times2}$, $h \in \R^2$.)
  \item Vodiš li egzaktan trag u razlomcima, uz decimale u zagradi?
\end{enumerate}
\end{predznanje}
